\begin{trapbox}
\textbf{\textcolor{CTrap}{易错一｜不判断是否存在两条实切线}}
\par\smallskip\begin{center}\begin{tikzpicture}[x=.88cm,y=.88cm,>=Stealth,font=\small]
\path[use as bounding box] (-1.8,-3.5) rectangle (13.2,4.2);
\foreach \shift/\px/\title/\value in {0/-1/两条实切线/4,5/0/一条实切线/0,10/1/无实切线/-4}{
\begin{scope}[shift={(\shift,0)}]
\node[font=\small\bfseries] at (.5,3.7) {\title};
\draw[->,gray] (-1.5,0)--(2.4,0) node[right] {$x$};
\draw[->,gray] (0,-2.8)--(0,2.9);
\draw[blue,thick,domain=-2.8:2.8,samples=61] plot ({\x*\x/4},\x);
\draw[red!75!black,very thick] ({-\px},-2.7)--({-\px},2.7);
\fill (\px,0) circle(2pt) node[below right,yshift=-6pt] {$P$};
\node[align=center] at (.4,-3.2) {$P(\px,0)$\\$y_0^2-4x_0=\value$};
\end{scope}}
\draw[orange!85!black,thick] (-1,0)--(1.6,2.6);
\draw[orange!85!black,thick] (-1,0)--(1.6,-2.6);
\fill (1,2) circle(2pt) node[below right] {$A$};
\fill (1,-2) circle(2pt) node[above right] {$B$};
\node[red!75!black,fill=white,inner sep=2pt] at (4,2.5) {$x=0$};
\node[red!75!black,fill=white,inner sep=2pt] at (8,2.5) {$x=-1$};
\end{tikzpicture}\end{center}\par\smallskip

上图固定 $y^2=4x$，只改变 $P$ 的位置。三幅红线均由公式得到，但只有左图红线连接两个不同切点；中图红线是顶点切线，右图红线与曲线无实交点。把三者都称为“切点弦”是错误的。
\par\textbf{纠正：} 对 $y^2=2px$，必须先查 $y_0^2-2px_0>0$；对 $x^2=2py$，必须先查 $x_0^2-2py_0>0$。等于零是\textbf{单条切线}，小于零\textbf{无实切点}。不能只说“外部”而不说明具体不等式。

\medskip
\textbf{\textcolor{CTrap}{易错二｜把 $P$ 当作切点弦上的点}}
切点弦经过两个\textbf{切点} $A,B$，一般不经过两切线的\textbf{交点} $P$。
\par\textbf{最小反例：} $y^2=4x$、$P(-1,0)$，切点弦为 $x=1$，而 $P$ 的横坐标为 $-1$。$P$ 明显不在弦上。
\par\textbf{进一步判断：} 若 $P$ 也满足 $y_0y=p(x+x_0)$，代入 $P$ 得 $y_0^2=2px_0$，这恰恰是\textbf{只有一条切线}的情形，而不是两个不同切点的情形。

\medskip
\textbf{\textcolor{CTrap}{易错三｜开口方向与公式混淆}}
\[
\begin{aligned}
 y^2=2px &:\quad \boxed{y_0y=p(x+x_0)},\\
 x^2=2py &:\quad \boxed{x_0x=p(y+y_0)}.
\end{aligned}
\]
\textbf{理解方法：} 哪个字母有平方，切点弦左边就由那个字母的两个坐标相乘，不用背一大段口诀。

\medskip
\textbf{\textcolor{CTrap}{易错四｜忘记系数减半或误写减号}}
把 $y_0y=p(x+x_0)$ 错写成 $y_0y=2p(x+x_0)$ 或 $y_0y=p(x-x_0)$。
\par\textbf{检查：} 若切点就是 $Q(u,v)$，切线 $vy=p(x+u)$ 代回 $Q$ 时右侧必须变成 $2pu$，与 $v^2=2pu$ 一致。

\medskip
\par\noindent\begin{minipage}{\linewidth}
\textbf{\textcolor{CTrap}{易错五｜将“直线有一个交点”一概叫作切线}}
\begin{center}\begin{tikzpicture}[x=1cm,y=1cm,>=Stealth,font=\small]
\path[use as bounding box] (-1.8,-4.2) rectangle (11,4.3);
\foreach \shift/\title in {0/水平线：横穿曲线,7/切线：与曲线方向一致}{
\begin{scope}[shift={(\shift,0)}]
\node[font=\small\bfseries] at (.6,3.8) {\title};
\draw[->,gray] (-1.4,0)--(3,0) node[right] {$x$};
\draw[->,gray] (0,-3)--(0,3.2);
\draw[blue,thick,domain=-2.9:3.1,samples=81] plot ({\x*\x/4},\x);
\fill (1,2) circle(2pt) node[below right] {$Q(1,2)$};
\end{scope}}
\draw[teal!80!black,thick] (-1.3,2)--(3,2);
\node[teal!80!black] at (-.5,2.6) {$y=2$};
\draw[orange!85!black,thick] (5.7,-.3)--(9.2,3.2);
\node[orange!85!black,anchor=west] at (8.4,1) {$y=x+1$};
\node[align=center] at (.7,-3.4) {只有一个交点，仍不是切线};
\node[align=center] at (7.7,-3.4) {重根对应切点，不发生横穿};
\end{tikzpicture}\end{center}\end{minipage}\par
判断切线时必须注意所用的方程形式。例如对 $y^2=2px$，采用 $x=x_0+m(y-y_0)$ 联立后是\textbf{关于 $y$ 的一元二次方程}，以判别式为零判相切；不能把水平直线 $y=c$ 与抛物线只有一个交点，就错误地当作抛物线切线。
\par\textbf{理由：} 如图，水平线虽只有一个实交点，却横穿曲线；右图切线与曲线在接触处方向一致。$y^2=2px$ 的真实切线都可写成 $vy=p(x+u)$（包括顶点处 $x=0$）。

\medskip
\textbf{\textcolor{CTrap}{易错六｜由弦方程反求 $P$ 时直接比较未经整理的系数}}
两直线方程可整体乘同一个非零数，故必须先把两边整理成同样的形式，再比较系数。
\par\textbf{例：} $y=x-1$ 与 $y_0y=2(x+x_0)$ 比较时，先写 $y=(2/y_0)x+2x_0/y_0$，得 $P(-1,2)$，而不是随意把 $y_0$ 读成 $1$。

\medskip
\textbf{\textcolor{CTrap}{易错七｜相加只得到中点关系，却误当作已求出中点坐标}}
两切点处切线代入 $P$ 后相加得到的是
\[y_0y_m=p(x_0+x_m),\]
这\textbf{只有一个等式、两个未知的中点坐标}，不能仅凭它就直接断言 $M$ 的坐标。
\par\textbf{纠正：} 用点斜式求弦时，只需这个关系即可消去中点坐标；如果题目要求 $M$ 的\textbf{具体坐标}，还应利用抛物线上两切点的方程相减，先得 $y_m=y_0$，再解得 $x_m=y_0^2/p-x_0$。

\medskip
\textbf{\textcolor{CTrap}{易错八｜默认 $k=p/y_0$ 永远可用，漏掉竖直切点弦}}
若 $y_0=0$，斜率式 $p/y_0$ 无意义，这不代表没有切点弦。
\par\textbf{最小反例：} $y^2=4x$、$P(-1,0)$，有两个切点 $(1,\pm 2)$，弦是 $x=1$，恰为竖直直线。
\par\textbf{纠正：} 保留没有除法的式子 $y_0(y_1-y_2)=p(x_1-x_2)$；当 $y_0=0$ 时，立即得到 $x_1=x_2=-x_0$，使用 $x=-x_0$，不再用点斜式。
\end{trapbox}
